% ============================================================
% Cover letter — submission to International Review of Economics & Finance
% Paper: "Momentum, Breadth, and Correlation in Tactical Asset Allocation:
%         A Rule-by-Rule Attribution"
% Authors: Garcia Martin I., Mases Campos M., Prior Sanz F. (corresponding)
% Standalone — compiles on its own (no bibliography needed).
% ============================================================
\documentclass[11pt,a4paper]{letter}
\usepackage[T1]{fontenc}
\usepackage[utf8]{inputenc}
\usepackage[margin=1in]{geometry}
\usepackage{setspace}
\usepackage{parskip}

\signature{Francesc Prior Sanz\\\textit{on behalf of all authors}}
\address{IQS School of Engineering\\Universitat Ramon Llull\\Via Augusta 390\\08017 Barcelona, Spain}

\begin{document}

\begin{letter}{The Editor\\\textit{International Review of Economics \& Finance}\\Elsevier}

\opening{Dear Editor,}

\onehalfspacing

We are pleased to submit our manuscript, \textbf{``Momentum, Breadth, and
Correlation in Tactical Asset Allocation: A Rule-by-Rule Attribution,''} for
consideration for publication in the \textit{International Review of Economics \& Finance}.

Rule-based tactical asset allocation (TAA) strategies are popular because they are
transparent: a few explicit rules (a momentum signal, a ranking cutoff, an
absolute-trend filter, a portfolio-construction criterion) drive the allocation,
rather than an opaque optimiser. Yet that very transparency exposes a question the
literature rarely answers directly: among the several rules a typical strategy
stacks, \emph{which} ones actually generate the out-of-sample edge, which are
inert, and is the result anything more than a repackaged trend premium? Without an
answer, a favourable backtest is hard to distinguish from the product of a
parameter search. The data-snooping literature typically delivers a single
corrected verdict on a \emph{strategy}; it less often asks \emph{which rule within
it} carries the edge. Our paper brings the two together.

We take one concrete cross-asset momentum strategy (a monthly rotation across 15
diversified ETFs, with no leverage, derivatives, or discretion) and treat it as a
case study. Every parameter is derived on a 2004--2015 training window and confirmed
on a 2016--2020 validation window, and the frozen strategy is evaluated once on a
locked 2021--2025 test block. We then decompose the out-of-sample performance
\emph{rule by rule}. The paper makes three contributions:

\begin{enumerate}\setlength{\itemsep}{2pt}
\item \textbf{A transparent, two-pronged attribution of the edge.} A matched
  ``skill-versus-luck'' permutation test replaces each rule's informed decision
  with a random draw from the same feasible set and finds that the momentum
  rule beats its matched random baseline on the Sharpe ratio ($p=0.0025$), while the
  minimum-correlation rule contributes through variance reduction ($p=0.010$) rather
  than in-sample return. A nested incremental spanning-alpha ladder then shows the edge
  arises from \emph{stacking} complementary rules, not from any single proprietary
  mechanism.
\item \textbf{The economic role of minimum-correlation construction.} We give the
  rule a formal optimality footing (it is the minimum-variance equal-weight
  portfolio under homogeneous screened returns) but show empirically that its value
  is \emph{risk control}, not return: it reduces variance and tail risk at
  statistically equal expected return. A head-to-head comparison against
  optimised constructions over the same eligible set shows that a long-only
  minimum-variance and an inverse-volatility risk-parity portfolio achieve
  comparable drawdowns but lower return and Sharpe, and neither beats the
  equal-weight minimum-correlation triplet, which captures the diversification
  benefit while avoiding the estimation error of covariance-matrix inversion.
\item \textbf{Robustness along every axis a referee is likely to probe.} The edge
  survives a significant spanning alpha against naive $1/N$ (locked test $+6.88\%$
  per annum, $p=0.008$); a residual alpha of $\approx +5\%$ per annum against
  joint controls for the equity, duration, and a generic trend (time-series
  momentum) factor built from the same assets, so it is not reducible to
  trend-following; realistic transaction costs; a Deflated Sharpe adjustment for
  the search; both sample halves; and a replication on 100\% real (non-proxy) ETF
  data ($+6.04\%$/yr, $p=0.0001$).
\end{enumerate}

We believe the paper fits the scope of the \textit{International Review of
Economics \& Finance}, which regularly publishes applied work on asset allocation,
momentum, and the empirical evaluation of investment strategies. Its methodological
contribution (a constructive, per-rule defence against the data-snooping
critique) should be of interest to readers beyond the specific strategy studied.

In the interest of a fair review, we are explicit about three limitations and how
we address each in the manuscript. \emph{First}, on \textbf{data snooping}: rather
than rest on a single corrected statistic, we localise the edge with a permutation
test and a nested spanning ladder (with bootstrap confidence intervals that make the
limited per-rule power explicit), and we report a Deflated Sharpe sensitivity table showing that the DSR
clears the 0.95 threshold at the actual trial count ($N=24$,
$\mathrm{DSR}=0.975$) and for any assumed count up to $N\approx154$; the primary defence against data snooping is nonetheless
the out-of-sample alpha and the matched-permutation attribution. \emph{Second}, on \textbf{transaction costs}: the strategy's
turnover is non-trivial ($\approx 10\times$ per year), so we evaluate net-of-cost
performance against the same cost schedule applied to all benchmarks and report a
break-even cost ($\approx 62$ bps one-way, roughly six times a realistic ETF cost),
together with a no-trade-band variant that lowers turnover while preserving the
alpha. \emph{Third}, on \textbf{benchmarks}: we state openly that the strategy does
\emph{not} beat a 60/40 portfolio on a risk-adjusted (Sharpe) basis; its
documented advantage is over naive $1/N$ and a Faber (2007) trend rule, and its
distinctive feature is materially lower drawdown (roughly half that of $1/N$,
60/40, and the S\&P~500). We prefer to make these points ourselves rather than
leave them for a referee to discover.

We confirm that this manuscript is original, has not been published previously, and
is not under consideration for publication elsewhere. All authors have approved the
submission and there are no conflicts of interest to declare. The strategy and all
empirical results derive from systematic, reproducible code; data-construction
choices and their limitations are documented in the manuscript.

Thank you for considering our work. We would be glad to respond to any questions and
look forward to the reviewers' feedback.

\closing{Sincerely,}

\vspace{1em}
\textit{Ignasi Garcia Martin \quad\textbullet\quad Marc Mases Campos
\quad\textbullet\quad Francesc Prior Sanz (corresponding author)}\\
IQS School of Engineering, Universitat Ramon Llull\\
Corresponding author: \texttt{francesc.prior@iqs.url.edu}

\end{letter}
\end{document}
